\documentclass[11pt]{article}
% Uses the canonical latex-docs semantic style in-repository. The local
% lightweight fallback is for independent compilation of this distribution.
\IfFileExists{technical-user-manual.sty}{\usepackage{technical-user-manual}}{\usepackage{../support/aaoa-standalone}}
\usepackage{mathtools}
\setDocTitle{Prerequisite-Based Learning Curriculum}
\setDocSubtitle{Sixteen weeks from proof skills to analytic engineering}
\setDocVersion{1.0}
\setDocStatus{Educational reference / unapproved working material}
\setDocOwner{Algorithm Engineering}
\setDocAuthor{Original companion material}
\setDocSummary{Sequential instruction, mastery gates, adaptive tracks and capstone performance evidence.}
\providecommand{\Expect}{\mathbb{E}}
\providecommand{\Prob}{\mathbb{P}}
\providecommand{\Var}{\operatorname{Var}}
\providecommand{\bigoh}{\mathcal{O}}
\begin{document}
\makedoctocpage

\section{Entry profile and diagnostic}
Designed for readers who can implement recursive algorithms and reason using discrete mathematics and elementary calculus. Before Week 1, demonstrate proof by induction, sums, logarithms, elementary probability (expectation and variance), simple tree traversals, and compiler/interpreter use. For a gap, complete a focused bridge: induction (2 sessions), combinatorics (2), probability (2), complex-number/power-series fundamentals (2).
\section{Sixteen-week prerequisite-led plan}
\begin{longtable}{p{0.09\linewidth}p{0.30\linewidth}p{0.44\linewidth}}\toprule Week & Reading / dependency & Required deliverable \\ \midrule
1 & Ch. 1.1--1.4; proof & Analysis-record template and diagnostic \\
2 & Ch. 1.5--1.8; expectation & Quicksort operation counter (PY01) \\
3 & Ch. 2.1--2.4; induction & WB01 first-order and characteristic roots \\
4 & Ch. 2.5--2.7; recursion & LAB-CPP01 mergesort \\
5 & Ch. 3.1--3.4; series algebra & WB02 coefficient derivations \\
6 & Ch. 3.5--3.11; probability & PGF and bivariate analysis \\
7 & Ch. 4.1--4.5; limits & WB03 error estimates \\
8 & Ch. 4.6--4.9; calculus & Mid-course assessment, 60-minute defense \\
9 & Ch. 5.1--5.3; classes & Catalan specification and OGF \\
10 & Ch. 5.4--5.5; singularities & Marked parameter and coefficient asymptotics \\
11 & Ch. 6.1--6.8; trees & PY03 BST shape experiments \\
12 & Ch. 6.9--6.15; complex trees & Tree-distribution comparison memo \\
13 & Ch. 7.1--7.9; permutations & CPP02 inversion-count validation \\
14 & Ch. 8.1--8.9; automata & PY04 trie/KMP analysis \\
15 & Ch. 9.1--9.8; occupancy & CPP02 chaining and coupon tests \\
16 & All previous methods & Capstone technical report and defense \\
\bottomrule\end{longtable}
\section{Weekly cadence}
Budget 6--9 focused hours: two hours for primary reading, two for derivations, two for implementation/testing, one for data analysis, and an additional one or two for corrective review. A week passes only when the original proof/check, executable artifact, and discrepancy note are complete. Weeks are sequencing targets, not claims about any learner's actual pace.
\section{Alternative paths}
\textbf{Analytic specialist:} weeks 1--10, then revisit Chapters 3--5 with bivariate generating functions and optional advanced singularity analysis. \textbf{Systems benchmarker:} weeks 1--4, 7--8, 11--16; carry independent operation models alongside all timing measurements. \textbf{Data structures specialist:} weeks 1--6, 9--16; verify Catalan and random BST are different laws.
\section{Competency matrix and thresholds}
\begin{longtable}{p{0.22\linewidth}p{0.50\linewidth}p{0.14\linewidth}}\toprule Domain & Evidence & Target \\ \midrule
Mathematical validity & Formally derived recurrence or coefficient identity with checked base cases & 80\% \\
Input-model literacy & Clearly named stochastic law and at least two limits & 80\% \\
Implementation accuracy & Exhaustive small-$n$ tests and reproducible seed logging & 90\% \\
Statistical reasoning & Raw samples, scale normalization and correct confidence claim & 80\% \\
Communication & Self-contained technical report with limitations and repeatable commands & 85\% \\
\bottomrule\end{longtable}
These thresholds are suite-defined learning targets, not official textbook grading criteria.
\section{Capstone options}
Choose one: (A) quicksort partition-strategy comparison with analytic and empirical models, (B) random-BST versus Catalan-tree structure study, or (C) trie versus hashing under varying prefix and occupancy distributions. Produce a proof annex, instrumented Python and C++ implementations or justified single-language variants, raw CSV, build manifest, error analysis, and a ten-minute defense answering \emph{which assumptions failed first?}

\end{document}
